\capitulo{2}{Objetivos del proyecto}

\section{Objetivo general}

El objetivo principal de este Trabajo Fin de Grado es diseñar e implementar un sistema SIEM ligero capaz de recoger, almacenar y analizar eventos de seguridad procedentes de diferentes dispositivos simulados dentro de un entorno de alta seguridad basado en un polvorín militar.

El sistema debe ser capaz de centralizar la información generada por distintos elementos de seguridad, aplicar mecanismos básicos de correlación de eventos y facilitar la detección de posibles incidentes mediante una interfaz de monitorización accesible y sencilla.

\section{Objetivos específicos}

Para alcanzar el objetivo general se plantean los siguientes objetivos específicos:

\begin{itemize}
\item Diseñar un entorno simulado de un polvorín militar que permita representar una infraestructura de alta seguridad.
\item Definir un formato de logs estructurado para facilitar el almacenamiento y tratamiento de la información generada por los distintos dispositivos.
\item Desarrollar un módulo de ingesta capaz de procesar los eventos generados por el entorno simulado.
\item Implementar una base de datos para almacenar de forma persistente los eventos y alertas detectados.
\item Desarrollar un motor de correlación basado en reglas que permita identificar situaciones potencialmente peligrosas.
\item Implementar mecanismos básicos de detección de anomalías sobre los eventos registrados.
\item Supervisar el estado de los sensores desplegados en el entorno monitorizado.
\item Desarrollar una API REST para consultar la información almacenada por el sistema.
\item Crear un dashboard web que permita visualizar eventos, alertas, anomalías y el estado de los dispositivos monitorizados.
\item Incorporar un sistema de notificaciones SMS simuladas para alertas críticas.
\item Diseñar una arquitectura modular que facilite futuras ampliaciones y el mantenimiento del sistema.
\item Evaluar la viabilidad de una solución SIEM ligera para entornos con recursos limitados.
\end{itemize}

\section{Objetivos técnicos}

Además de los objetivos funcionales anteriores, durante el desarrollo del proyecto se han perseguido una serie de objetivos técnicos:

\begin{itemize}
\item Aplicar conocimientos adquiridos durante el grado relacionados con programación, bases de datos y desarrollo de aplicaciones.
\item Utilizar tecnologías actuales de desarrollo como Python, SQLite, FastAPI y Streamlit.
\item Diseñar una solución modular y fácilmente mantenible.
\item Aplicar buenas prácticas de desarrollo y control de versiones mediante Git y GitHub.
\item Obtener experiencia práctica en el desarrollo de sistemas relacionados con la monitorización y gestión de eventos de seguridad.
\end{itemize}
